\documentclass[border=3pt]{standalone}
% 公共导言区 —— 所有 TikZ 图 \input 本文件后使用
\usepackage[UTF8, fontset=mac]{ctex}
\usepackage{tikz}
\usetikzlibrary{arrows.meta, positioning, calc, shapes.geometric, decorations.pathmorphing, backgrounds, fit}
\definecolor{zhusha}{HTML}{9B2D20}
\definecolor{mohei}{HTML}{1A1A1A}
\definecolor{shenhui}{HTML}{555555}
\definecolor{qianhui}{HTML}{F5F0EC}
\definecolor{lan}{HTML}{1F4E79}
\definecolor{qing}{HTML}{2E7D6E}
\tikzset{
  block/.style={draw=shenhui, fill=white, thick, rounded corners=1.5pt,
                font=\small, align=center, inner sep=4pt, minimum height=7mm},
  core/.style={block, fill=lan!12, draw=lan, font=\small\bfseries},
  periph/.style={block, fill=qing!10, draw=qing},
  power/.style={block, fill=zhusha!8, draw=zhusha},
  note/.style={font=\footnotesize, text=shenhui, align=center},
  lab/.style={font=\footnotesize, text=mohei},
  sig/.style={-{Stealth[length=2.2mm]}, thick, color=mohei},
  fbsig/.style={-{Stealth[length=2.2mm]}, thick, color=zhusha, dashed},
  vec/.style={-{Stealth[length=3mm]}, very thick, color=zhusha},
}

\begin{document}
\begin{tikzpicture}[x=1mm, y=1mm, font=\small, >=Stealth]

% ===== 三相轴 A/B/C（互差 120°）=====
\foreach \ang/\n in {0/A, 120/B, 240/C} {
  \draw[->, thick, shenhui] (0,0) -- (\ang:34mm);
  \node[lab, text=shenhui] at (\ang:37mm) {$\n$ 相轴};
}
% ===== αβ 轴 =====
\draw[->, very thick, mohei] (0,0) -- (0:44mm) node[right, lab] {$\alpha$};
\draw[->, very thick, mohei] (0,0) -- (90:44mm) node[above, lab] {$\beta$};
% ===== dq 轴（随转子旋转 θe）=====
\draw[->, very thick, zhusha] (0,0) -- (38:46mm) node[right, lab, text=zhusha] {$d$（转子磁场方向）};
\draw[->, very thick, zhusha] (0,0) -- (128:46mm) node[above, lab, text=zhusha] {$q$（垂直，出力方向）};

% ===== 电流矢量 =====
\draw[vec] (0,0) -- (75:40mm) node[above right, lab, text=zhusha] {$\vec{i}_s$（定子电流矢量）};

% ===== 角度标注 =====
\draw[thin, mohei] (14mm, 0) arc (0:38:14mm);
\node[lab] at (19:17mm) {$\theta_e$};
\draw[thin, zhusha] (10mm, 0) arc (0:75:10mm);
\node[lab, text=zhusha] at (40:12.5mm) {$\gamma$};
\node[note, text=zhusha] at (52:15mm) {FOC 目标：$\gamma = 90°$};

% ===== 旋转方向 =====
\draw[->, thick, qing] (38mm, -6mm) arc (0:-210:7mm);
\node[note, text=qing] at (24mm, -14mm) {$\omega$（电角速度）};

% ===== 相电流投影小弧 =====
\node[note, anchor=north west, align=left, text=shenhui] at (-46mm, -34mm) {
  abc：三相静止系（物理绕组轴）\\
  $\alpha\beta$：两相静止系（Clarke 变换）\\
  dq：两相旋转系（Park 变换，随转子以 $\omega$ 旋转）\\
  同一个 $\vec{i}_s$，在 dq 系里是两个直流量 $i_d$、$i_q$};

% 电流矢量在 dq 轴上的投影虚线
\coordinate (tip) at (75:40mm);
\coordinate (dend) at (38:46mm);
\coordinate (qend) at (128:46mm);
\draw[dashed, thin, zhusha] (tip) -- ($(0,0)!(tip)!(dend)$) node[midway, above left=-0.5mm, note, text=zhusha] {$i_q$};
\draw[dashed, thin, zhusha] (tip) -- ($(0,0)!(tip)!(qend)$);

\end{tikzpicture}
\end{document}
